\cvname{林彦}

\cvtagline{助理教授 / 计算机科学系 / 奥尔堡大学}

电子邮箱：\href{mailto:lyan@cs.aau.dk}{lyan@cs.aau.dk} \,/\,
个人主页：\href{https://www.yanlincs.com}{yanlincs.com} \,/\,
电话：+45 9186-1833

{\color{rule}\rule{\linewidth}{0.4pt}}

\vspace{\entryskip}
我是一名计算机科学领域的研究者，专注于数据挖掘，即用机器学习模型从大规模数据中发现并利用模式。
目前的研究方向包括轨迹表示学习、时空数据挖掘和人工智能材料科学。
此外，我正持续积累教学经验，并将科研背景与个人兴趣驱动的计算机科学探索融入教学。

\vspace{\entryskip}
我已在 SIGMOD、KDD、NeurIPS、AAAI 和 IEEE TKDE 等顶级会议和期刊发表论文 33 篇，其中 17 篇以第一作者、共同第一作者或主要学生作者身份完成。
我在 \href{https://scholar.google.com/citations?user=nHMmG2UAAAAJ}{谷歌学术} 上的总被引次数超过 1,000，h 指数为 15，i10 指数为 23。


\section*{教育经历}

\roleentry{北京交通大学}{计算机科学博士}{2019 年 9 月至 2024 年 3 月}
{博士论文：时空轨迹数据的自监督学习。导师：万怀宇教授。}

% \roleentry{奥尔堡大学}{计算机科学访问博士生}{2022 年 9 月至 2023 年 9 月}
% {导师：Bin Yang 教授和 Jilin Hu 教授。}

\roleentry{北京交通大学}{计算机科学学士}{2015 年 9 月至 2019 年 6 月}{}


\section*{工作经历}

\roleentry{奥尔堡大学计算机科学系}{助理教授}{2026 年 6 月至今}
{}

\roleentry{奥尔堡大学计算机科学系}{博士后研究员}{2024 年 6 月至 2026 年 5 月}
{与奥尔堡大学玻璃结构与力学研究组合作开展人工智能材料科学研究。}
